% Part: lambda-calculus
% Chapter: lambda-definability
% Section: primitive-recursive-functions

\documentclass[../../../include/open-logic-section]{subfiles}

\begin{document}

\olfileid[fr]{lam}{rep}{prf}
\olsection{Fonctions récursives primitives \texorpdfstring{$\lambda$}{λ}-définissables}

Rappelons que les fonctions récursives
primitives s'obtiennent à partir des
fonctions initiales $\Zero$, $\Succ$ et
$\Proj{n}{i}$ par composition et par
récurrence primitive.

\begin{lem}
\ollabel{lem:basic}
Les fonctions initiales $\Zero$,
$\Succ$ et les projections
$\Proj{n}{i}$ sont $\lambda$-définissables.
\end{lem}

\begin{proof}
Les termes suivants les définissent :
\begin{align*}
  \fn{Zero} & \ident
    \lambd[a][\lambd[fx][x]]\\
  \fn{Succ} & \ident
    \lambd[a][\lambd[fx][f(a f x)]]\\
  \fn{Proj}^n_i & \ident
    \lambd[x_0\dots x_{n-1}][x_i].
\end{align*}
\end{proof}

\begin{lem}
\ollabel{lem:comp}
Supposons que $f$ soit une fonction
totale à $k$ arguments et que
$g_0,\dots,g_{k-1}$ soient des fonctions
totales à $n$ arguments, toutes
$\lambda$-définissables respectivement
par $F$ et $G_0,\dots,G_{k-1}$.
La fonction $h$ obtenue par composition
est alors $\lambda$-définissable.
\end{lem}

\begin{proof}
Un terme qui définit~$h$ est
\[
  H \ident
  \lambd[x_0\dots x_{n-1}]
    [F\,(G_0 x_0\dots x_{n-1})
       \dots(G_{k-1}x_0\dots x_{n-1})].
\]
La vérification est laissée en exercice.
\end{proof}

\begin{prob}
Compléter la preuve du
\olref[lam][rep][prf]{lem:comp} : montrer
que, pour tous $m_0,\dots,m_{n-1}$,
$H\,\num{m_0}\dots\num{m_{n-1}}
\red\num{h(m_0,\dots,m_{n-1})}$.
\end{prob}

Le \olref{lem:comp} n'exige pas
que $f$ ou les $g_i$ soient récursives
primitives. Pour l'argument donné ici,
leur totalité et leur
$\lambda$-définissabilité suffisent.

\begin{lem}\ollabel{lem:prim}
Supposons que $f$ soit une fonction
totale à $n$ arguments et~$g$ une
fonction totale à $n+2$ arguments,
$\lambda$-définies par $F$ et~$G$.
Si $h$ est obtenue de $f$ et~$g$
par récurrence primitive, alors $h$
est $\lambda$-définissable.
\end{lem}

\begin{proof}
La fonction~$h$ est déterminée par
\begin{align*}
  h(x_1,\dots,x_n,0)
    &=f(x_1,\dots,x_n),\\
  h(x_1,\dots,x_n,y+1)
    &=g(x_1,\dots,x_n,y,
        h(x_1,\dots,x_n,y)).
\end{align*}
La définition itère une transition
qui emploie l'indice courant et la
valeur déjà obtenue. Un numéral de
Church peut compter ces itérations.

Pour alléger l'écriture, traitons
un seul paramètre~$x$ en plus de
l'argument de récurrence~$y$.
Notons $D_x$ la transition sur les
couples et posons
\begin{align*}
  D_x &\ident
    \lambd[p][\tuple{
      \fn{Succ}(\fn{Fst}\,p),
      G\,x\,(\fn{Fst}\,p)\,(\fn{Snd}\,p)}],\\
  H &\ident
    \lambd[x][\lambd[y][
      \fn{Snd}(y D_x
      \tuple{\num{0},F\,x})]].
\end{align*}
L'état est un couple dont la première
composante est l'indice courant et la
seconde la valeur correspondante de~$h$.
Chaque application de $D_x$ augmente
l'indice et calcule la valeur suivante;
à la fin, $\fn{Snd}$ renvoie cette valeur.

Fixons $n$ et écrivons $D_n$ pour le
terme obtenu de $D_x$ en remplaçant
$x$ par~$\num{n}$. Montrons par
récurrence sur~$m$ que
\[
  D_n^m\tuple{\num{0},F\,\num{n}}
  \red\tuple{\num{m},\num{h(n,m)}}.
\]
Pour $m=0$, le terme de gauche est
$\tuple{\num{0},F\,\num{n}}$.
Puisque $F$ définit~$f$, il se réduit
à $\tuple{\num{0},\num{f(n)}}$,
soit $\tuple{\num{0},\num{h(n,0)}}$.

Supposons maintenant l'énoncé vrai
pour~$m$. L'application de $D_n$
au couple obtenu donne
\begin{align*}
  D_n^{m+1}\tuple{\num{0},F\,\num{n}}
    &\red D_n\,
      \tuple{\num{m},\num{h(n,m)}}\\
    &\red
      \tuple{\fn{Succ}(\num{m}),
        G\,\num{n}\,\num{m}\,
          \num{h(n,m)}}\\
    &\red
      \tuple{\num{m+1},
        \num{g(n,m,h(n,m))}}\\
    &\ident
      \tuple{\num{m+1},
        \num{h(n,m+1)}}.
\end{align*}
Le troisième pas utilise les termes
qui définissent $\Succ$ et~$g$;
la dernière égalité est l'équation
de récurrence. L'invariant est établi.

Enfin, sur les entrées~$\num{n}$ et
$\num{m}$, le terme~$H$ donne
\begin{align*}
  H\,\num{n}\,\num{m}
    &\red \fn{Snd}
      (\num{m}D_n
        \tuple{\num{0},F\,\num{n}})\\
    &\red \fn{Snd}
      (D_n^m\tuple{\num{0},F\,\num{n}})\\
    &\red \fn{Snd}
      \tuple{\num{m},\num{h(n,m)}}\\
    &\red \num{h(n,m)}.
\end{align*}
Il définit donc~$h$ dans ce cas;
les autres paramètres se traitent
de la même manière.
\end{proof}

\begin{prop}
Toute fonction récursive primitive
est $\lambda$-définissable.
\end{prop}

\begin{proof}
Les fonctions initiales sont
$\lambda$-définissables par le
\olref{lem:basic}. Le
\olref{lem:comp} et le \olref{lem:prim}
donnent la fermeture par composition
et par récurrence primitive.
\end{proof}

À chaque étape de la récurrence, le couple fournit le
numéral du résultat précédent. Dans une récurrence partielle stricte,
un résultat précédent indéfini doit rendre la nouvelle
valeur indéfinie. Un terme représentant $g$ peut ignorer
l'argument qui lui transmet ce résultat et produire une
forme normale. La construction de ce lemme garantit donc
la récurrence des fonctions totales.

\end{document}
